%====================================
% VORLÄUFIGE GLIEDERUNG
%====================================
\section{Vorläufige Gliederung der Bachelor-Thesis}

Die nachfolgende Tabelle zeigt die vorläufige Gliederung im Überblick.

\begin{table}[H]
\centering
\begin{tabularx}{\textwidth}{Xp{2cm}}
\toprule
\textbf{Kapitel / Abschnitt} & \textbf{ca. Seiten} \\
\midrule
Verzeichnisse (Inhalt, Abkürzungen, Abb., Tab.) & i\,--\,v \\
\midrule
1. Einleitung & 3\,--\,4 \\
\quad 1.1 Problemstellung und Relevanz & \\
\quad 1.2 Zielsetzung und Forschungsfragen & \\
\quad 1.3 Aufbau der Arbeit & \\
\midrule
2. Theoretische Grundlagen & 12\,--\,13 \\
\quad 2.1 Enterprise-Resource-Planning-Systeme & \\
\quad 2.2 Erfolg von ERP-Implementierungen (DeLone \& McLean) & \\
\quad 2.3 Change Management (ADKAR-Modell) & \\
\quad 2.4 Vorläufiger Bezugsrahmen der Untersuchung & \\
\midrule
3. Methodisches Vorgehen & 5\,--\,6 \\
\quad 3.1 Systematische Literaturanalyse (Webster \& Watson) & \\
\quad 3.2 Suchstrategie, Datenbanken und Schlagworte & \\
\quad 3.3 Auswahlprozess und Auswertungslogik (Mayring) & \\
\midrule
4. Ergebnisse der Literaturanalyse & 10\,--\,12 \\
\quad 4.1 Deskriptive Auswertung des Samples & \\
\quad 4.2 Identifizierte CM-Faktoren (Kommunikation, Führung, Schulung) & \\
\quad 4.3 Verknüpfung mit DeLone \& McLean-Dimensionen & \\
\midrule
5. Handlungsempfehlungen für den Mittelstand & 5--6 \\
\quad 5.1 ADKAR-basierte Empfehlungen & \\
\quad 5.2 Kritische Würdigung & \\
\midrule
6. Fazit und Ausblick & 3\,--\,4 \\
\quad 6.1 Beantwortung der Forschungsfragen & \\
\quad 6.2 Limitationen (Publication Bias, Heterogenit\"at des Mittelstands) \newline und Implikationen & \\
\midrule
Literaturverzeichnis & \\
\midrule
\textbf{Gesamtumfang Textteil} & \textbf{38\,--\,45} \\
\bottomrule
\end{tabularx}
\caption{Vorläufige Gliederung der Bachelor-Thesis}
\label{tab:gliederung}
\end{table}
